\noindent\textit{Proof.}
Let \(\mathcal O\) be the finite support of \(O\), put
\(J=|\mathcal O|\), and write
\[
v=\bigl(P_{\mathrm{obs}}(O=o)\bigr)_{o\in\mathcal O}\in\mathbb R^J,
\qquad
\widehat v_n=
\bigl(\widehat P_n(O=o)\bigr)_{o\in\mathcal O}.
\]
Assumption \(\mathrm{ass{:}iid\text{-}sampling}\) and Lemma
\(\mathrm{lem{:}finite\text{-}multinomial\text{-}central\text{-}limit}\)
give
\[
\sqrt n(\widehat v_n-v)\rightsquigarrow
\mathbb Z_{P_{\mathrm{obs}}}
\quad\text{in }\mathbb R^J.
\tag{1}
\]

We first construct the reduced-support map and use finite-dimensional
polyhedral directional calculus to derive its derivative.  Set
\(\mathbb D=\mathbb R^J\), \(\mathbb E=\mathbb R^2\), and
\(\mathbb D_0=\mathbb D\).  For every \(w\in\mathbb D\), compute the
ratio analogues prescribed by Definition
\(\mathrm{def{:}plugin\text{-}endpoint\text{-}estimator}\), using zero
when a conditioning denominator is zero.  Form the raw capacity contrasts,
apply \(\Pi_{\mathfrak C}\), and compute
\(q_0,q_1,\Delta q,m,B_L,B_U\).  Sum numerator and denominator
contributions only over the fixed set \(\mathcal X_+(P)\), divide by the
common denominator when it is nonzero, and otherwise use \((0,1)\).
Denote this total map by
\[
\Phi:\mathbb D_\phi=\mathbb D\longrightarrow\mathbb E.
\tag{2}
\]
The values assigned at zero denominators merely totalize \(\Phi\) away
from the population neighborhood.

If \(x\in\mathcal X_+(P)\), then \(p_xm(x)>0\), so \(p_x>0\).
Assumption \(\mathrm{ass{:}instrument\text{-}overlap}\) gives
\[
P_{\mathrm{obs}}(X=x,Z=1)=p_x\pi(x)
\ge p_x\varepsilon_Z>0
\]
and
\[
P_{\mathrm{obs}}(X=x,Z=0)=p_x\{1-\pi(x)\}
\ge p_x\varepsilon_Z>0.
\tag{3}
\]
Thus every ratio denominator used by \(\Phi\) on the fixed support is
positive at \(v\) and remains nonzero in a neighborhood of \(v\).  The
map from \(w\) to the finite vector of cell probabilities and raw
capacities is continuously differentiable there by the quotient rule.

At the population law, Proposition
\(\mathrm{prop{:}capacity\text{-}identification}\), which is part of the
proof basis of Theorem
\(\mathrm{thm{:}sharp\text{-}exact\text{-}mass\text{-}threshold\text{-}interval}\),
shows that the raw capacity vector is the nonnegative vector
\(c=(\ell,h)\).  For a raw-capacity direction \(g\), coordinatewise
positive part has Hadamard directional derivative
\[
[\Pi'_{\mathfrak C,c}(g)]_k=
\begin{cases}
g_k,&c_k>0,\\
\max\{g_k,0\},&c_k=0.
\end{cases}
\tag{4}
\]
Indeed, for every \(t_r\downarrow0\) and \(g_r\to g\), the corresponding
scalar difference quotient converges to the displayed value; finiteness
of the coordinates makes the convergence Euclidean.  Let
\(\dot\ell_i(x),\dot h_j(x)\) denote the directions after (4), and set
\[
\dot q_0(x)=\sum_i\dot\ell_i(x),\qquad
\dot q_1(x)=\sum_j\dot h_j(x),\qquad
\dot{\Delta q}(x)=\dot q_1(x)-\dot q_0(x).
\tag{5}
\]

For the lower endpoint, define
\[
A_{x,t}=\ell_{\le t}(x)-h_{\le t}(x),\qquad
G_x=\max_{t\in\mathcal Y}A_{x,t},\qquad
s_x=\min\{\Delta q(x),0\},
\tag{6}
\]
and the nonempty active set
\[
\mathcal A_L(x)=\{t\in\mathcal Y:A_{x,t}=G_x\}.
\]
Writing
\[
\dot A_{x,t}=\sum_{i\le t}\dot\ell_i(x)
-\sum_{j\le t}\dot h_j(x),
\]
the finite active-set expansion gives
\[
\dot G_x=\max_{t\in\mathcal A_L(x)}\dot A_{x,t}.
\tag{7}
\]
To verify (7) along arbitrary convergent directions, note that every
inactive threshold lies a strictly positive distance below \(G_x\); the
smallest such gap is positive because \(\mathcal Y\) is finite.  No
inactive threshold can therefore maximize for sufficiently small step
size, and division by that step leaves the maximum of the active
first-order terms.  Direct scalar case analysis gives
\[
\dot s_x=
\begin{cases}
\dot{\Delta q}(x),&\Delta q(x)<0,\\
0,&\Delta q(x)>0,\\
\min\{\dot{\Delta q}(x),0\},&\Delta q(x)=0.
\end{cases}
\tag{8}
\]
Since \(B_L(x)=\max\{0,G_x+s_x\}\), the same scalar expansion gives
\[
\dot B_L(x)=
\begin{cases}
\dot G_x+\dot s_x,&G_x+s_x>0,\\
0,&G_x+s_x<0,\\
\max\{\dot G_x+\dot s_x,0\},&G_x+s_x=0.
\end{cases}
\tag{9}
\]
Equations (7)--(9) allow arbitrary simultaneous threshold ties.  In
particular, on \(\Delta q(x)=0\), equation (8) is exactly the asserted
face derivative \(g\mapsto\min\{\dot{\Delta q}(x),0\}\).

For the upper endpoint, define
\[
U_{x,t}=\ell_{<t}(x)+h_{>t}(x),
\qquad
\dot U_{x,t}=\sum_{i<t}\dot\ell_i(x)+\sum_{j>t}\dot h_j(x).
\tag{10}
\]
The survivor-mass derivative, obtained by the scalar minimum case split,
is
\[
\dot m(x)=
\begin{cases}
\dot q_0(x),&q_0(x)<q_1(x),\\
\dot q_1(x),&q_1(x)<q_0(x),\\
\min\{\dot q_0(x),\dot q_1(x)\},&q_0(x)=q_1(x).
\end{cases}
\tag{11}
\]
Index the mass cap by \(\star\), put
\(V_{x,\star}=m(x)\), \(\dot V_{x,\star}=\dot m(x)\), and put
\(V_{x,t}=U_{x,t}\), \(\dot V_{x,t}=\dot U_{x,t}\) for
\(t\in\mathcal Y\).  With
\[
\mathcal A_U(x)=
\{a\in\{\star\}\cup\mathcal Y:V_{x,a}=B_U(x)\},
\]
the same finite-gap argument, now for a minimum, yields
\[
\dot B_U(x)=\min_{a\in\mathcal A_U(x)}\dot V_{x,a}.
\tag{12}
\]
Thus (11)--(12) include positive-survivor zero-gap cells and arbitrary
ties between the mass cap and any number of threshold cuts.

The sequential arguments above establish Hadamard directional
differentiability: every expansion holds for each \(t_r\downarrow0\) and
each convergent sequence of directions.  The derivative maps are finite
maxima or minima of continuous linear maps and coordinatewise positive
parts, hence are continuous in the direction.  If \(\dot p_x\) denotes
the cell-probability direction, define, for \(e=L,U\),
\[
N_e=\sum_{x\in\mathcal X_+(P)}p_xB_e(x),
\qquad
\dot N_e=\sum_{x\in\mathcal X_+(P)}
\{\dot p_xB_e(x)+p_x\dot B_e(x)\},
\tag{13}
\]
and
\[
M=\sum_{x\in\mathcal X_+(P)}p_xm(x),
\qquad
\dot M=\sum_{x\in\mathcal X_+(P)}
\{\dot p_xm(x)+p_x\dot m(x)\}.
\tag{14}
\]
Cells outside \(\mathcal X_+(P)\) have \(p_xm(x)=0\), so (14) equals
the aggregate survivor mass in the frozen core.  Assumption
\(\mathrm{ass{:}positive\text{-}aggregate\text{-}survivors}\) gives
\(M>0\), which is the premise needed for division.  The quotient
derivative is therefore
\[
\left(\frac{N_e}{M}\right)'(g)
=\frac{\dot N_eM-N_e\dot M}{M^2},
\qquad e\in\{L,U\}.
\tag{15}
\]
Equations (3)--(15) prove that \(\Phi\) is Hadamard directionally
differentiable at \(v\), tangentially to
\(\mathbb D_0=\mathbb R^J\), with the displayed continuous derivative.

We now apply the directional delta method.  The spaces
\(\mathbb D=\mathbb R^J\) and \(\mathbb E=\mathbb R^2\) are Banach
spaces; \(\mathbb D_\phi=\mathbb D\); the base point is \(v\); the
estimator is \(\widehat v_n\); and the rate is \(\sqrt n\to\infty\).
Weak convergence is (1).  The finite-dimensional Gaussian vector is
tight and supported in \(\mathbb R^J=\mathbb D_0\).  Lemma
\(\mathrm{lem{:}hadamard\text{-}directional\text{-}delta\text{-}method}\)
therefore gives
\[
\sqrt n\{\Phi(\widehat v_n)-\Phi(v)\}
=\Phi'_v\!\left(\sqrt n(\widehat v_n-v)\right)+o_P(1)
\rightsquigarrow\Phi'_v(\mathbb Z_{P_{\mathrm{obs}}}).
\tag{16}
\]

It remains to compare the fixed-support map with the screened estimator.
Write \(r_x=p_xm(x)\).  Since \(M>0\), the finite set
\(\mathcal X_+(P)\) is nonempty and
\[
\rho=\min_{x\in\mathcal X_+(P)}r_x>0.
\tag{17}
\]
Equation (1) implies \(\widehat v_n\to_Pv\).  For every
\(x\in\mathcal X_+(P)\), (3) and continuity of the local ratio,
projection, and finite-minimum maps give \(\widehat r_x\to_Pr_x\).
Ordinary convergence \(\eta_n\to0\) implies that eventually
\(0<\eta_n<\rho/2\).  Finiteness of \(\mathcal X_+(P)\) then shows that
all its cells are retained with probability tending to one; no
monotonicity of \(n\mapsto\eta_n\) is used.

If \(p_x=0\), independent sampling gives \(\widehat p_x=0\) almost
surely, so \(\widehat r_x=0\).  Suppose instead that \(p_x>0\) and
\(m(x)=0\).  Nonnegativity of \(q_0(x)\) and \(q_1(x)\) implies that at
least one total, say \(q_a(x)\), is zero.  Every constituent population
capacity in that arm is then zero.  Overlap and \(p_x>0\) make the
relevant conditioning denominators positive, so the corresponding raw
empirical contrasts are smooth near \(v\), centered at zero, and are
\(O_P(n^{-1/2})\) by (1).  Coordinatewise positive part preserves this
order.  Since \(K\) is fixed and \(0\le\widehat p_x\le1\),
\[
0\le\widehat m(x)\le\widehat q_a(x)=O_P(n^{-1/2}),
\qquad
0\le\widehat r_x=\widehat p_x\widehat m(x)
\le\widehat m(x)=O_P(n^{-1/2})=o_P(\eta_n),
\tag{18}
\]
where the last equality follows from
\(\sqrt n\eta_n\to\infty\).  A finite union over cells combines
(17)--(18) into
\[
P\!\left\{\{x:\widehat r_x>\eta_n\}=\mathcal X_+(P)\right\}
\longrightarrow1.
\tag{19}
\]

On the event in (19), the screened sums in Definition
\(\mathrm{def{:}plugin\text{-}endpoint\text{-}estimator}\) are exactly
the fixed-support sums defining \(\Phi(\widehat v_n)\).  Their common
denominator is positive because the retained set is nonempty and each
retained \(\widehat r_x\) exceeds the positive number \(\eta_n\); hence
the fallback is not used.  At the population law, Definitions
\(\mathrm{def{:}threshold\text{-}cuts}\) and
\(\mathrm{def{:}identified\text{-}interval}\), together with Theorem
\(\mathrm{thm{:}sharp\text{-}exact\text{-}mass\text{-}threshold\text{-}interval}\),
give \(\Phi(v)=\Psi(P_{\mathrm{obs}})\).  Replacing
\(\Phi(\widehat v_n)\) by \(\widehat\Psi_n\) in (16) on an event whose
probability tends to one proves the asserted directional limit.  The same
replacement and (16) yield
\(\widehat\Psi_n\to_P\Psi(P_{\mathrm{obs}})\).  Thus positivity,
ordinary convergence \(\eta_n\to0\), and
\(\sqrt n\eta_n\to\infty\) suffice for all three conclusions; no
direction separation, monotonicity of the screening sequence,
multiplier calibration, or uniform-in-law conclusion is used.
\(\square\)